\chapter{Conclusion}
\label{ch:conclusion}

This Major Project-I report has developed, rigorously evaluated, and experimentally validated an end-to-end automated robotic vision inspection system for precision surface defect detection on thermoformed industrial packaging casings. By replacing conventional horizontal bounding boxes with an orientation-aware YOLOv8-OBB architecture and identifying the destructive impact of spatial mosaic augmentation on micro-defects, this research resolves longstanding feature dilution bottlenecks in automated optical quality assurance.

%------------------------------------------------------------------
\section{Achievement of Research Objectives}
\label{sec:achievement}

The project has successfully fulfilled all five targeted research and engineering objectives formulated in Chapter~\ref{ch:intro}:

\begin{enumerate}
\item \textbf{Achievement of Objective 1 (Dataset Engineering):} Compiled and validated an exhaustive industrial benchmark dataset of 5,065 high-resolution images annotated across nine defect classes with ground-truth oriented bounding polygons $[x, y, w, h, \theta]$, capturing 24,477 total defect instances.
\item \textbf{Achievement of Objective 2 (Optomechanical Station):} Designed and fabricated an industrial optomechanical imaging enclosure integrating coaxial diffuse ring LED lighting, telecentric optics, and high-speed optical triggering that completely eliminates ambient lighting variance and specular reflections.
\item \textbf{Achievement of Objective 3 (Model Formulation):} Formulated and implemented an optimized single-stage YOLOv8-OBB model incorporating decoupled orientation heads, Gaussian Wasserstein Distance (GWD) loss, Varifocal classification loss (VFL), and Distribution Focal Loss (DFL), eliminating angular boundary discontinuities.
\item \textbf{Achievement of Objective 4 (Augmentation Optimization):} Discovered and empirically demonstrated that deactivating spatial Mosaic augmentation and constraining geometric scaling prevents the collapse of sub-pixel scratch features, elevating micro-scratch mAP50 by an extraordinary +83.24 percentage points.
\item \textbf{Achievement of Objective 5 (Edge Deployment \& Actuation):} Serialized and deployed the trained model via ONNX-Slim and TensorRT FP16 on an industrial edge computer with ROS2 middleware and PLC synchronization, achieving a deterministic 24.0~ms cycle latency (41.6~FPS) and enabling real-time pneumatic ejection.
\end{enumerate}

%------------------------------------------------------------------
\section*{Key Project Outcomes}
\addcontentsline{toc}{section}{Key Project Outcomes}

The empirical investigations and factory validation yielded the following landmark quantitative outcomes:

\begin{itemize}
\item \textbf{Landmark mAP50 Accuracy of 97.55\%:} The optimized YOLOv8-OBB model achieved an overall mean Average Precision at IoU 0.50 (mAP50) of 97.55\% across all nine defect categories, outperforming the baseline generic model (49.02\%) by +48.53 percentage points.
\item \textbf{Micro-Scratch Detection Rate Elevated to 95.85\%:} Detection mAP50 on slender micro-scratch abrasions surged from 12.61\% to 95.85\% (+83.24\%), resolving the most critical defect escape vulnerability in packaging manufacturing.
\item \textbf{Deterministic 24.0~ms Real-Time Edge Latency:} The edge-optimized TensorRT inference pipeline delivers a deterministic end-to-end processing latency of 24.0~ms (41.6~FPS), comfortably satisfying the 25.0~ms line cycle threshold for 120 units/minute conveyor sorting.
\item \textbf{Defect Escape Rate Reduced to 0.3\%:} The factory defect escape rate plummeted from 14.5\% (manual inspection) and 8.2\% (horizontal AOI) to an ultra-low 0.3\%, virtually eliminating commercial product recall risks.
\item \textbf{Rapid Economic Payback of 4.2 Months:} Economic modeling demonstrates annual net operational savings of \$540,000 for a 63M units/year facility, achieving full capital investment payback within 4.2 months.
\end{itemize}

%------------------------------------------------------------------
\section{Future Scope}
\label{sec:future}

While the current system achieves industry-leading performance on visible surface anomalies, several promising research avenues will be pursued in subsequent development phases:

\begin{enumerate}
\item \textbf{Multi-Spectral \& Hyperspectral SWIR Imaging:} Integrating Short-Wave Infrared (SWIR, 900--1700~nm) or hyperspectral illumination to inspect multi-layer barrier films (such as EVOH core layers) for internal delaminations and trapped moisture pockets invisible to visible-spectrum optics.
\item \textbf{INT8 / FP8 Edge Quantization:} Implementing post-training INT8 and FP8 quantization via TensorRT with calibration-aware training to compress model size to under 10~MB and reduce inference latency below 12~ms, enabling deployment on ultra-low-power micro-edge microcontrollers.
\item \textbf{6-DoF Robotic Gripper Integration:} Expanding the downstream rejection station from simple pneumatic blow-off nozzles to a high-speed 6-DoF delta robot arm equipped with adaptive soft-pneumatic grippers, enabling gentle extraction of delicate or heavy packaging containers without deformation~\cite{zhang2024robotic}.
\item \textbf{Cross-Domain Transfer Learning to Medical Packaging:} Applying few-shot transfer learning and domain adaptation techniques to generalize the trained YOLOv8-OBB backbone to pharmaceutical blister packs, medical syringe trays, and semiconductor carrier tapes with minimal retraining.
\end{enumerate}
